% Chunk: \S4.2, Creation and Annihilation Operators.
% Source: Weinberg QFT Vol. I, printed pp. 173--177 (PDF pp. 196--200),
% ending immediately before the \S4.3 heading.
% Status: source-reviewed, compile-clean, and visually checked; frozen.

\subsection{Creation and Annihilation Operators}

Creation and annihilation operators may be defined in terms of their effect
on the normalized multi-particle states discussed in the previous section.
The \textit{creation operator} $a^\dagger_q$ (or in more detail, $a^\dagger_{\mathbf p,\sigma,n}$) is defined as the
operator that simply adds a particle with quantum numbers $q$ at the front
of the list of particles in the state
\begin{equation}
  a^\dagger_q\ket{q_1q_2\cdots q_N}=\ket{q q_1q_2\cdots q_N}.
  \tag{4.2.1}
  \label{eq:4.2.1}
\end{equation}
In particular, the $N$-particle state can be obtained by acting on the vacuum
with $N$ creation operators
\begin{equation}
  a^\dagger_{q_1}a^\dagger_{q_2}\cdots a^\dagger_{q_N}\ket{\Omega}
  =\ket{q_1\cdots q_N}.
  \tag{4.2.2}
  \label{eq:4.2.2}
\end{equation}
It is conventional for this operator to be called $a^\dagger_q$; its adjoint, which
is then called $a_q$, may be calculated from Eq.~\eqref{eq:4.1.7}. As we shall
now show, $a_q$ removes a particle from any state on which it acts, and
is therefore known as an \textit{annihilation operator}. In particular, when the
particles $q,q_1,\ldots,q_N$ are either all bosons or all fermions, we have
\begin{equation}
  a_q\ket{q_1q_2\cdots q_N}
  =\sum_{r=1}^{N}(\pm)^{r+1}\delta(q-q_r)
    \ket{q_1\cdots q_{r-1}q_{r+1}\cdots q_N},
  \tag{4.2.3}
  \label{eq:4.2.3}
\end{equation}
with a $+1$ or $-1$ sign for bosons or fermions, respectively. (Here is the
proof. We want to calculate the scalar product of $a_q\ket{q_1q_2\cdots q_N}$ with an
arbitrary state $\ket{q'_1\cdots q'_M}$. Using Eq.~\eqref{eq:4.2.1}, this is
\[
  \bra{q'_1\cdots q'_M}a_q\ket{q_1\cdots q_N}
  =\bigl(a^\dagger_q\ket{q'_1\cdots q'_M}\bigr)^\dagger
    \ket{q_1\cdots q_N}
  =\braket{q q'_1\cdots q'_M}{q_1\cdots q_N}.
\]
We now use Eq.~\eqref{eq:4.1.7}. The sum over permutations $\mathcal P$ of $1,2,\ldots,N$
can be written as a sum over the integer $r$ that is permuted into the
first place, i.e. $\mathcal P r=1$, and over mappings $\overline{\mathcal P}$ of the remaining integers
$1,\ldots,r-1,r+1,\ldots,N$ into $1,\ldots,N-1$. Furthermore, the sign factor is
\[
  \delta_{\mathcal P}=(\pm)^{r-1}\delta_{\overline{\mathcal P}}
\]
with upper and lower signs for bosons and fermions, respectively. Hence,
using Eq.~\eqref{eq:4.1.7} twice,
\begin{align*}
  \bra{q'_1\cdots q'_M}a_q\ket{q_1\cdots q_N}
  ={}& \delta_{N,M+1}
  \sum_{r=1}^{N}\sum_{\overline{\mathcal P}}(\pm)^{r-1}
  \delta_{\overline{\mathcal P}}\delta(q-q_r)
  \prod_{i=1}^{M}\delta(q'_i-q_{\overline{\mathcal P}i})
  \\
  ={}& \delta_{N,M+1}\sum_{r=1}^{N}(\pm)^{r-1}
  \braket{q'_1\cdots q'_M}{q_1\cdots q_{r-1}q_{r+1}\cdots q_N}.
\end{align*}
Both sides of Eq.~\eqref{eq:4.2.3} thus have the same matrix element with any state
$\ket{q'_1\cdots q'_M}$, and are therefore equal, as was to be shown.) As a special case
of Eq.~\eqref{eq:4.2.3}, we note that for both bosons and fermions, $a_q$ annihilates
the vacuum
\begin{equation}
  a_q\ket{\Omega}=0.
  \tag{4.2.4}
  \label{eq:4.2.4}
\end{equation}

As defined here, the creation and annihilation operators satisfy an important
commutation or anticommutation relation. Applying the operator
$a_{q'}$ to Eq.~\eqref{eq:4.2.1} and using Eq.~\eqref{eq:4.2.3} gives
\begin{align*}
  a_{q'}a^\dagger_q\ket{q_1\cdots q_N}
  ={}& \delta(q'-q)\ket{q_1\cdots q_N}
  \\
  &+\sum_{r=1}^{N}(\pm)^{r+2}\delta(q'-q_r)
    \ket{q q_1\cdots q_{r-1}q_{r+1}\cdots q_N}.
\end{align*}
(The sign in the second term is $(\pm)^{r+2}$ because $q_r$ is in the $(r+1)$-th place
in $\ket{q q_1\cdots q_N}$.) On the other hand, applying the operator $a^\dagger_q$ to Eq.~\eqref{eq:4.2.3}
gives
\[
  a^\dagger_q a_{q'}\ket{q_1\cdots q_N}
  =\sum_{r=1}^{N}(\pm)^{r+1}\delta(q'-q_r)
  \ket{q q_1\cdots q_{r-1}q_{r+1}\cdots q_N}.
\]
Subtracting or adding, we have then
\[
  \bigl[a_{q'}a^\dagger_q\mp a^\dagger_q a_{q'}\bigr]\ket{q_1\cdots q_N}
  =\delta(q'-q)\ket{q_1\cdots q_N}.
\]
But this holds for all states $\ket{q_1\cdots q_N}$ (and may easily be seen to hold also
for states containing both bosons and fermions) and therefore implies the
operator relation
\begin{equation}
  a_{q'}a^\dagger_q\mp a^\dagger_q a_{q'}=\delta(q'-q).
  \tag{4.2.5}
  \label{eq:4.2.5}
\end{equation}
In addition, Eq.~\eqref{eq:4.2.2} gives immediately
\begin{equation}
  a^\dagger_{q'}a^\dagger_q\mp a^\dagger_q a^\dagger_{q'}=0
  \tag{4.2.6}
  \label{eq:4.2.6}
\end{equation}
and so also
\begin{equation}
  a_{q'}a_q\mp a_q a_{q'}=0.
  \tag{4.2.7}
  \label{eq:4.2.7}
\end{equation}
As always, the top and bottom signs apply for bosons and fermions,
respectively. According to the conventions discussed in the previous
section, the creation and/or annihilation operators for particles of two
different species commute if either particle is a boson, and anticommute
if both are fermions.

The above discussion could have been presented in reverse order (and
in most textbooks usually is). That is, we could have started with the
commutation or anticommutation relations Eqs.~\eqref{eq:4.2.5}--\eqref{eq:4.2.7}, derived
from the canonical quantization of some given field theory. Multi-particle
states would have then been defined by Eq.~\eqref{eq:4.2.2}, and their scalar
products Eq.~\eqref{eq:4.1.7} derived from the commutation or anticommutation
relations. In fact, as discussed in Chapter 1, such a treatment would be
much closer to the way that this formalism developed historically. We
have followed an unhistorical approach here because we want to free
ourselves from any dependence on pre-existing field theories, and rather
wish to understand why field theories are the way they are.

We will now prove the fundamental theorem quoted at the beginning
of this chapter: \textit{any operator $\mathcal O$ may be expressed as a sum of products of
creation and annihilation operators}
\begin{align}
  \mathcal O={}&\sum_{N=0}^{\infty}\sum_{M=0}^{\infty}
  \int \dd q'_1\cdots \dd q'_N\,\dd q_1\cdots \dd q_M
  \notag\\
  &\quad\cdot a^\dagger_{q'_1}\cdots a^\dagger_{q'_N}
  a_{q_M}\cdots a_{q_1}
  \notag\\
  &\quad\cdot C_{NM}(q'_1\cdots q'_N q_1\cdots q_M).
  \tag{4.2.8}
  \label{eq:4.2.8}
\end{align}
That is, we want to show that the $C_{NM}$ coefficients can be chosen to give
the matrix elements of this expression any desired values. We do this by
mathematical induction. First, it is trivial that by choosing $C_{00}$ properly,
we can give $\bra{\Omega}\mathcal O\ket{\Omega}$ any desired value, irrespective of the values of $C_{NM}$
with $N>0$ and/or $M>0$. We need only use Eq.~\eqref{eq:4.2.4} to see that
Eq.~\eqref{eq:4.2.8} has the vacuum expectation value
\[
  \bra{\Omega}\mathcal O\ket{\Omega}=C_{00}.
\]
Now suppose that the same is true for all matrix elements of $\mathcal O$ between
$N$- and $M$-particle states, with $N<L,M\leq K$ or $N\leq L,M<K$; that
is, that these matrix elements have been given some desired values by
an appropriate choice of the corresponding coefficients $C_{NM}$. To see that
the same is then also true of matrix elements of $\mathcal O$ between any $L$- and
$K$-particle states, use Eq.~\eqref{eq:4.2.8} to evaluate
\begin{align*}
  \bra{q'_1\cdots q'_L}\mathcal O\ket{q_1\cdots q_K}
  ={}& L!K!C_{LK}(q'_1\cdots q'_L q_1\cdots q_K)
  \\
  &+\text{terms involving $C_{NM}$ with $N<L,M\leq K$ or $N\leq L,M<K$.}
\end{align*}
Whatever values have already been given to $C_{NM}$ with $N<L,M\leq K$
or $N\leq L,M<K$, there is clearly some choice of $C_{LK}$ which gives this
matrix element any desired value.

Of course, an operator need not be expressed in the form~\eqref{eq:4.2.8}, with
all creation operators to the left of all annihilation operators. (This
is often called the `normal' order of the operators.) However, if the
formula for some operator has the creation and annihilation operators
in some other order, we can always bring the creation operators to the
left of the annihilation operators by repeated use of the commutation or
anticommutation relations, picking up new terms from the delta function
in Eq.~\eqref{eq:4.2.5}.

For instance, consider any sort of additive operator $F$ (like momentum,
charge, etc.) for which
\begin{equation}
  F\ket{q_1\cdots q_N}
  =\bigl(f(q_1)+\cdots+f(q_N)\bigr)\ket{q_1\cdots q_N}.
  \tag{4.2.9}
  \label{eq:4.2.9}
\end{equation}
Such an operator can be written as in Eq.~\eqref{eq:4.2.8}, but using only the term
with $N=M=1$:
\begin{equation}
  F=\int \dd q\,a^\dagger_q a_q f(q).
  \tag{4.2.10}
  \label{eq:4.2.10}
\end{equation}
In particular, the free-particle Hamiltonian is always
\begin{equation}
  H_0=\int \dd q\,a^\dagger_q a_q E(q)
  \tag{4.2.11}
  \label{eq:4.2.11}
\end{equation}
where $E(q)$ is the single-particle energy
\[
  E(\mathbf p,\sigma,n)=\sqrt{\mathbf p^2+m_n^2}.
\]

We will need the transformation properties of the creation and annihilation
operators for various symmetries. First, let's consider inhomogeneous
proper orthochronous Lorentz transformations. Recall that the $N$-particle
states have the Lorentz transformation property
\begin{align*}
  &U_0(\Lambda,a)\ket{\mathbf p_1\sigma_1n_1;\,\mathbf p_2\sigma_2n_2;\,\ldots}
  \\
  &\quad={}
  e^{-\ii(\Lambda p_1)\cdot a}e^{-\ii(\Lambda p_2)\cdot a}\cdots
  \sqrt{\frac{(\Lambda p_1)^0(\Lambda p_2)^0\cdots}{p_1^0p_2^0\cdots}}
  \notag\\
  &\qquad\cdot\sum_{\bar\sigma_1\bar\sigma_2\cdots}
  D_{\bar\sigma_1\sigma_1}^{(j_1)}\bigl(W(\Lambda,p_1)\bigr)
  D_{\bar\sigma_2\sigma_2}^{(j_2)}\bigl(W(\Lambda,p_2)\bigr)\cdots
  \notag\\
  &\qquad\cdot
  \ket{\mathbf p_{1\Lambda}\bar\sigma_1n_1;\,
       \mathbf p_{2\Lambda}\bar\sigma_2n_2;\,\ldots}.
\end{align*}
Here $\mathbf p_\Lambda$ is the three-vector part of $\Lambda p$, $D_{\bar\sigma\sigma}^{(j)}(R)$ is the same unitary spin-$j$
representation of the three-dimensional rotation group as used in Section
2.5, and $W(\Lambda,p)$ is the particular rotation
\[
  W(\Lambda,p)=L^{-1}(\Lambda p)\Lambda L(p),
\]
where $L(p)$ is the standard `boost' that takes a particle of mass $m$ from rest
to four-momentum $p^\mu$. (Of course, $m$ and $j$ depend on the species label $n$.
This is all for $m\ne0$; we will return to the massless particle case in the
following chapter.) Now, these states can be expressed as in Eq.~\eqref{eq:4.2.2}
\[
  \ket{\mathbf p_1\sigma_1n_1;\,\mathbf p_2\sigma_2n_2;\,\ldots}
  =a^\dagger_{\mathbf p_1,\sigma_1,n_1}a^\dagger_{\mathbf p_2,\sigma_2,n_2}\cdots
  \ket{\Omega},
\]
where $\ket{\Omega}$ is the Lorentz-invariant vacuum state
\[
  U_0(\Lambda,a)\ket{\Omega}=\ket{\Omega}.
\]
In order that the state~\eqref{eq:4.2.2} should transform properly, it is necessary
and sufficient that the creation operator have the transformation rule
\begin{align}
  &U_0(\Lambda,a)a^\dagger_{\mathbf p,\sigma,n}U_0^{-1}(\Lambda,a)
  \notag\\
  &\qquad={}
  e^{-\ii(\Lambda p)\cdot a}\sqrt{\frac{(\Lambda p)^0}{p^0}}
  \sum_{\bar\sigma}D_{\bar\sigma\sigma}^{(j)}\bigl(W(\Lambda,p)\bigr)
  a^\dagger_{\mathbf p_\Lambda,\bar\sigma,n}.
  \tag{4.2.12}
  \label{eq:4.2.12}
\end{align}
In the same way, the operators $C$, $P$, and $T$, that induce charge-conjugation,
space inversion, and time-reversal transformations on free particle states\footnote{We
omit the subscript `0' on these operators, because in virtually all cases where $C$,
$P$, and/or $T$ are conserved, the operators that induce these transformations on
`in' and `out' states are the same as those defined by their action on free-particle
states. This is not the case for continuous Lorentz transformations, for which it is
necessary to distinguish between the operators $U(\Lambda,a)$ and $U_0(\Lambda,a)$.}
transform the creation operators as:
\begin{equation}
  Ca^\dagger_{\mathbf p,\sigma,n}C^{-1}
  =\xi_n a^\dagger_{\mathbf p,\sigma,n^c},
  \tag{4.2.13}
  \label{eq:4.2.13}
\end{equation}
\begin{equation}
  Pa^\dagger_{\mathbf p,\sigma,n}P^{-1}
  =\eta_n a^\dagger_{-\mathbf p,\sigma,n},
  \tag{4.2.14}
  \label{eq:4.2.14}
\end{equation}
\begin{equation}
  Ta^\dagger_{\mathbf p,\sigma,n}T^{-1}
  =\zeta_n(-1)^{j-\sigma}a^\dagger_{-\mathbf p,-\sigma,n}.
  \tag{4.2.15}
  \label{eq:4.2.15}
\end{equation}

As mentioned in the previous section, although we have been dealing
with operators that create and annihilate particles in free-particle states,
the whole formalism can be applied to `in' and `out' states, in which case
we would introduce operators $a_{\mathrm{in}}$ and $a_{\mathrm{out}}$ defined in the same way by
their action on these states. These operators satisfy a Lorentz transformation
rule just like Eq.~\eqref{eq:4.2.12}, but with the true Lorentz transformation
operator $U(\Lambda,a)$ instead of the free-particle operator $U_0(\Lambda,a)$.
